% ID: MAT-TRI-EQG-017
% Titolo: Dominio e zeri di una funzione goniometrica (1)
% Disciplina: Matematica
% Area: trigonometria
% Argomento: equazioni_goniometriche
% Classe: Quarta secondaria di secondo grado
% Difficolta: 3
% Tipo: problema_numerico
% Tag: matematica, trigonometria, equazioni_goniometriche
% Fonte: verifica_goniometria_base_anonimizzata
% Provenienza: materiale fornito dal docente; testo anonimizzato; soluzione redatta per Exergo

\begin{esercizio}
Determina il dominio e gli zeri della funzione: 
\[
f(x)=\dfrac{\sin x-\dfrac{\sqrt{2}}{2}}{1+\tan x}.
\]
Salvo diversa indicazione, lavora in $\mathbb R$; angoli in radianti.
\end{esercizio}

\begin{soluzione}
Il dominio esclude $x=\pi/2+k\pi$ (tangente non definita)
e $x=-\pi/4+k\pi$ (denominatore nullo).
Il numeratore si annulla per $x=\pi/4+2k\pi$ e $3\pi/4+2k\pi$;
solo la prima famiglia appartiene al dominio. In tutte le famiglie di soluzioni $k\in\mathbb{Z}$.
\end{soluzione}
